\documentclass[11pt,a4paper]{article}\usepackage[]{graphicx}\usepackage[]{xcolor}
% maxwidth is the original width if it is less than linewidth
% otherwise use linewidth (to make sure the graphics do not exceed the margin)
\makeatletter
\def\maxwidth{ %
  \ifdim\Gin@nat@width>\linewidth
    \linewidth
  \else
    \Gin@nat@width
  \fi
}
\makeatother

\definecolor{fgcolor}{rgb}{0.345, 0.345, 0.345}
\newcommand{\hlnum}[1]{\textcolor[rgb]{0.686,0.059,0.569}{#1}}%
\newcommand{\hlstr}[1]{\textcolor[rgb]{0.192,0.494,0.8}{#1}}%
\newcommand{\hlcom}[1]{\textcolor[rgb]{0.678,0.584,0.686}{\textit{#1}}}%
\newcommand{\hlopt}[1]{\textcolor[rgb]{0,0,0}{#1}}%
\newcommand{\hlstd}[1]{\textcolor[rgb]{0.345,0.345,0.345}{#1}}%
\newcommand{\hlkwa}[1]{\textcolor[rgb]{0.161,0.373,0.58}{\textbf{#1}}}%
\newcommand{\hlkwb}[1]{\textcolor[rgb]{0.69,0.353,0.396}{#1}}%
\newcommand{\hlkwc}[1]{\textcolor[rgb]{0.333,0.667,0.333}{#1}}%
\newcommand{\hlkwd}[1]{\textcolor[rgb]{0.737,0.353,0.396}{\textbf{#1}}}%
\let\hlipl\hlkwb

\usepackage{framed}
\makeatletter
\newenvironment{kframe}{%
 \def\at@end@of@kframe{}%
 \ifinner\ifhmode%
  \def\at@end@of@kframe{\end{minipage}}%
  \begin{minipage}{\columnwidth}%
 \fi\fi%
 \def\FrameCommand##1{\hskip\@totalleftmargin \hskip-\fboxsep
 \colorbox{shadecolor}{##1}\hskip-\fboxsep
     % There is no \\@totalrightmargin, so:
     \hskip-\linewidth \hskip-\@totalleftmargin \hskip\columnwidth}%
 \MakeFramed {\advance\hsize-\width
   \@totalleftmargin\z@ \linewidth\hsize
   \@setminipage}}%
 {\par\unskip\endMakeFramed%
 \at@end@of@kframe}
\makeatother

\definecolor{shadecolor}{rgb}{.97, .97, .97}
\definecolor{messagecolor}{rgb}{0, 0, 0}
\definecolor{warningcolor}{rgb}{1, 0, 1}
\definecolor{errorcolor}{rgb}{1, 0, 0}
\newenvironment{knitrout}{}{} % an empty environment to be redefined in TeX

\usepackage{alltt}
\usepackage[utf8]{inputenc}
\usepackage[T1]{fontenc}
\usepackage[margin=1in]{geometry}
\usepackage{graphicx}
\usepackage{fancyhdr}
\usepackage{booktabs}
\usepackage{array}
\usepackage{longtable}
\usepackage{enumitem}
\usepackage{xcolor}
\usepackage{titlesec}
\usepackage{hyperref}
\usepackage{amsmath}

% Header and footer
\pagestyle{fancy}
\fancyhf{}
\fancyhead[L]{SOP-TOC-SSM-001}
\fancyhead[R]{Revision 1.0}
\fancyfoot[C]{\thepage}
\renewcommand{\headrulewidth}{0.4pt}
\renewcommand{\footrulewidth}{0.4pt}

% Section formatting
\titleformat{\section}
  {\normalfont\Large\bfseries}{\thesection}{1em}{}
\titleformat{\subsection}
  {\normalfont\large\bfseries}{\thesubsection}{1em}{}
\titleformat{\subsubsection}
  {\normalfont\normalsize\bfseries}{\thesubsubsection}{1em}{}

% Hyperref setup
\hypersetup{
    colorlinks=true,
    linkcolor=blue,
    filecolor=magenta,      
    urlcolor=cyan,
    pdftitle={SOP: Shimadzu TOC Analyzer with SSM},
    pdfpagemode=FullScreen,
}
\IfFileExists{upquote.sty}{\usepackage{upquote}}{}
\begin{document}

% Title page
\begin{titlepage}
    \centering
    \vspace*{2cm}
    
    {\Huge\bfseries Standard Operating Procedure\par}
    \vspace{1cm}
    {\Large Shimadzu SSM-5000A TOC Analyzer\\for Soil Samples with Varying Carbonate Concentrations\par}
    \vspace{2cm}
    
    \begin{tabular}{ll}
        \toprule
        \textbf{Document Number:} & SOP-TOC-SSM-001 \\
        \textbf{Revision:} & 1.0 \\
        \textbf{Effective Date:} & October 2025 \\
        \textbf{Review Date:} & October 2026 \\
        \bottomrule
    \end{tabular}
    
    \vfill
    
    {\large \textbf{Document Control:} This is a controlled document. Printed copies are uncontrolled and may not be current.\par}
    
\end{titlepage}

% Table of contents
\tableofcontents
\newpage

\section{Purpose and Scope}

This SOP describes the procedure for determining Total Carbon (TC), Total Inorganic Carbon (TIC), and Total Organic Carbon (TOC) in soil samples using the Shimadzu SSM-5000A TOC Analyzer equipped with a Solid Sample Module. This procedure is applicable to soil samples with high ($>10\%$), medium (2--10\%), and low ($<2\%$) carbonate concentrations.

\section{Principle}

The Shimadzu SSM-5000A system uses high-temperature combustion (900°C) to convert all carbon in the sample to CO$_2$, which is then quantified by NDIR (Non-Dispersive Infrared) detection. For carbonate-bearing soils:

\begin{itemize}[noitemsep]
    \item \textbf{TC (Total Carbon)} = measured by direct combustion
    \item \textbf{TIC (Total Inorganic Carbon)} = measured by acid treatment (200°C)
    \item \textbf{TOC (Total Organic Carbon)} = TC $-$ TIC
\end{itemize}

\section{Safety and Precautions}

\subsection{Personal Protective Equipment (PPE)}
\begin{itemize}[noitemsep]
    \item Safety glasses or goggles
    \item Laboratory coat
    \item Heat-resistant gloves when handling hot ceramic boats
    \item Nitrile gloves when handling acids
\end{itemize}

\subsection{Chemical Hazards}
\begin{itemize}[noitemsep]
    \item Phosphoric acid (25\% solution): corrosive, causes severe burns
    \item High-purity oxygen: oxidizer, keep away from combustibles
    \item Soil samples: may contain pathogens or heavy metals
\end{itemize}

\subsection{Equipment Hazards}
\begin{itemize}[noitemsep]
    \item Combustion furnace reaches 900°C
    \item Acid furnace reaches 200°C
    \item Allow adequate cooling time before handling boats
\end{itemize}

\section{Equipment and Materials}

\subsection{Instrumentation}
\begin{itemize}[noitemsep]
    \item Shimadzu SSM-5000A TOC Analyzer with Solid Sample Module
    \item High-purity oxygen supply (99.999\% minimum purity)
    \item Carrier gas supply (oxygen, 99.999\% minimum purity)
    \item Computer with SSM-5000A Control Software (v3.x or higher)
    \item Analytical balance (0.0001 g precision, capacity $\geq$ 220 g)
    \item Autosampler unit (40-position carousel)
\end{itemize}

\subsection{Consumables}
\begin{itemize}[noitemsep]
    \item Ceramic sample boats (platinum or high-purity alumina, pre-combusted)
    \item 25\% phosphoric acid solution (prepared fresh weekly)
    \item Glucose anhydrous (certified reference material, 99.9\% purity)
    \item Calcium carbonate standard (reagent grade, dried at 105°C for 2 hours)
    \item Quartz wool for combustion tube packing
    \item Halogen moisture trap cartridge
\end{itemize}

\subsection{Sample Preparation Equipment}
\begin{itemize}[noitemsep]
    \item Mortar and pestle (agate or alumina)
    \item Ball mill (optional, for fine grinding)
    \item 2 mm and 0.25 mm sieves (stainless steel or brass)
    \item Drying oven (105°C $\pm$ 2°C)
    \item Desiccator with indicating desiccant (drierite or silica gel)
\end{itemize}

\section{Reagents and Standards}

\subsection{Phosphoric Acid Solution (25\% v/v)}
\begin{enumerate}[noitemsep]
    \item Add 250 mL concentrated phosphoric acid (85\%) slowly to 600 mL deionized water in a 1000 mL volumetric flask
    \item \textbf{CAUTION:} Add acid to water, never water to acid
    \item Mix thoroughly by swirling (exothermic reaction)
    \item Allow to cool to room temperature
    \item Dilute to 1000 mL with deionized water
    \item Store in amber glass bottle with screw cap
    \item Label with: name, concentration, preparation date, expiration date (7 days)
    \item Prepare fresh weekly or when contamination is suspected
\end{enumerate}

\subsection{TOC Calibration Standards (from Glucose)}
\textbf{Stock Solution Preparation (1000 mg C/L):}
\begin{enumerate}[noitemsep]
    \item Weigh exactly 2.500 g $\pm$ 0.001 g anhydrous glucose into 1000 mL volumetric flask
    \item Glucose contains 40.0\% carbon by weight
    \item Dissolve in approximately 800 mL deionized water
    \item Dilute to volume with deionized water
    \item Mix thoroughly by inversion (20 times)
    \item Store at 4°C for up to 7 days
\end{enumerate}

\textbf{Working Standard Preparation:}

\begin{table}[h]
\centering
\small
\begin{tabular}{ccccc}
\toprule
\textbf{Level} & \textbf{Target Conc.} & \textbf{Stock Vol.} & \textbf{Final Vol.} & \textbf{Mass C} \\
 & \textbf{(mg C/L)} & \textbf{(mL)} & \textbf{(mL)} & \textbf{(mg)} \\
\midrule
1 & 10 & 1.0 & 100 & 0.001 \\
2 & 25 & 2.5 & 100 & 0.0025 \\
3 & 50 & 5.0 & 100 & 0.005 \\
4 & 100 & 10.0 & 100 & 0.010 \\
5 & 250 & 25.0 & 100 & 0.025 \\
6 & 500 & 50.0 & 100 & 0.050 \\
\bottomrule
\end{tabular}
\caption{TC Calibration Standard Preparation from Stock Solution}
\end{table}

\subsection{TIC Calibration Standards (from Calcium Carbonate)}

\textbf{Preparation:}
\begin{enumerate}[noitemsep]
    \item Dry reagent-grade CaCO$_3$ at 105°C for 2 hours
    \item Cool in desiccator for minimum 30 minutes
    \item Calculate theoretical carbon content: CaCO$_3$ molecular weight = 100.09 g/mol; C = 12.01 g/mol
    \item Carbon content = $\frac{12.01}{100.09} \times 100 = 12.0\%$ by weight
    \item Prepare solid standards by weighing directly into ceramic boats:
\end{enumerate}

\begin{table}[h]
\centering
\begin{tabular}{ccccc}
\toprule
\textbf{Level} & \textbf{CaCO$_3$ Mass} & \textbf{Carbon Mass} & \textbf{Expected Peak} & \textbf{Acid Volume} \\
 & \textbf{(mg)} & \textbf{(mg C)} & \textbf{Area (est.)} & \textbf{($\mu$L)} \\
\midrule
1 & 8.3 & 1.0 & Low & 150 \\
2 & 41.7 & 5.0 & Medium-Low & 150 \\
3 & 83.3 & 10.0 & Medium & 200 \\
4 & 208.3 & 25.0 & Medium-High & 200 \\
5 & 416.7 & 50.0 & High & 250 \\
\bottomrule
\end{tabular}
\caption{TIC Calibration Standard Preparation}
\end{table}

\begin{itemize}[noitemsep]
    \item Store prepared standards in desiccator
    \item Use within 24 hours of preparation
    \item Weigh standards to nearest 0.1 mg
\end{itemize}

\section{Sample Preparation}

\subsection{Initial Processing}
\begin{enumerate}[noitemsep]
    \item Air-dry soil samples at room temperature (3--7 days) or in drying oven at 40°C maximum (to prevent organic matter loss)
    \item Remove large debris, roots, stones, and organic fragments manually
    \item Break up aggregates gently using mortar and pestle
    \item Pass through 2 mm sieve for initial homogenization
    \item Record percent of sample retained on sieve (coarse fraction)
\end{enumerate}

\subsection{Fine Grinding}
\begin{enumerate}[noitemsep]
    \item Take representative subsample (approximately 50 g) from $<$2 mm fraction
    \item Further grind using agate mortar and pestle or ball mill
    \item Grind until all material passes through 0.25 mm sieve (60 mesh)
    \item For ball mill: grind for 5--10 minutes at medium speed
    \item Mix thoroughly by quartering method or tumbling for 2 minutes
    \item Store ground samples in airtight containers (polypropylene or glass)
\end{enumerate}

\subsection{Moisture Determination}
\begin{enumerate}[noitemsep]
    \item Weigh 2--5 g of ground sample into pre-weighed, pre-dried aluminum crucible
    \item Record crucible + wet sample weight to 0.0001 g
    \item Dry at 105°C $\pm$ 2°C for 24 hours in forced-air oven
    \item Cool in desiccator for minimum 30 minutes
    \item Reweigh crucible + dry sample to 0.0001 g
    \item Calculate moisture content:
\end{enumerate}

\begin{equation}
\text{Moisture \%} = \frac{\text{Wet weight} - \text{Dry weight}}{\text{Wet weight}} \times 100
\end{equation}

\begin{itemize}[noitemsep]
    \item Perform moisture determination in triplicate
    \item Use average moisture content for all calculations
    \item Redetermine moisture if sample storage exceeds 7 days
\end{itemize}

\subsection{Sample Weighing Guidelines}

\begin{longtable}{p{4cm}p{3cm}p{3cm}p{3cm}}
\toprule
\textbf{Soil Type} & \textbf{TC Analysis} & \textbf{TIC Analysis} & \textbf{Expected TIC} \\
\midrule
\endfirsthead
\multicolumn{4}{c}%
{\tablename\ \thetable\ -- \textit{Continued from previous page}} \\
\toprule
\textbf{Soil Type} & \textbf{TC Analysis} & \textbf{TIC Analysis} & \textbf{Expected TIC} \\
\midrule
\endhead
\midrule
\multicolumn{4}{r}{\textit{Continued on next page}} \\
\endfoot
\bottomrule
\endlastfoot
High Carbonate ($>10\%$) & 10--20 mg & 5--10 mg & $>1.2\%$ C \\
Medium Carbonate (2--10\%) & 20--40 mg & 10--30 mg & 0.24--1.2\% C \\
Low Carbonate ($<2\%$) & 40--80 mg & 30--80 mg & $<0.24\%$ C \\
Non-calcareous & 50--100 mg & N/A & 0\% C \\
\end{longtable}

\section{Detailed Instrument Startup Procedure}

\subsection{Pre-Startup Inspection (Day Before Analysis)}

\subsubsection{Gas Supply Verification}
\begin{enumerate}[noitemsep]
    \item Locate oxygen cylinder and check pressure gauge
    \item Verify cylinder pressure $>$ 500 psi (34 bar)
    \item If pressure $<$ 500 psi, arrange for cylinder replacement
    \item Inspect regulator for damage, corrosion, or leaks
    \item Verify gas line connections are tight (hand-tight + 1/4 turn with wrench)
    \item Check for leaks using leak detection solution or soap solution:
    \begin{itemize}
        \item Apply solution to all fittings and connections
        \item Open cylinder valve slightly
        \item Look for bubble formation indicating leaks
        \item Tighten connections if bubbles appear
        \item Retest after tightening
    \end{itemize}
    \item Verify gas tubing is not kinked, cracked, or discolored
\end{enumerate}

\subsubsection{Instrument Physical Inspection}
\begin{enumerate}[noitemsep]
    \item Open instrument doors and visually inspect:
    \begin{itemize}
        \item Combustion tube for cracks or deposits
        \item Ceramic boats for contamination or damage
        \item Quartz wool packing (replace if discolored or compressed)
        \item Acid delivery system for crystallization or leaks
        \item NDIR detector cell window for fogging or deposits
    \end{itemize}
    \item Check autosampler carousel:
    \begin{itemize}
        \item Verify smooth rotation (rotate manually)
        \item Inspect positioning mechanism for alignment
        \item Check boat holder clips for wear
    \end{itemize}
    \item Verify drain bottles are empty (empty if $>$ 50\% full)
    \item Check desiccant color in moisture trap (replace if pink/saturated)
\end{enumerate}

\subsection{System Power-Up Sequence (Analysis Day)}

\subsubsection{Step 1: Turn On Gas Supply}
\begin{enumerate}[noitemsep]
    \item Slowly open oxygen cylinder main valve (turn counterclockwise)
    \item Turn valve 1.5--2 full rotations (not fully open)
    \item Observe high-pressure gauge on cylinder (should read cylinder pressure)
    \item Adjust regulator to set delivery pressure:
    \begin{itemize}
        \item Turn regulator adjustment knob clockwise to increase pressure
        \item Set to 0.4--0.5 MPa (58--73 psi) as indicated on low-pressure gauge
        \item For SSM-5000A: recommended setting is 0.45 MPa (65 psi)
    \end{itemize}
    \item Verify stable pressure reading (no fluctuations)
    \item Record cylinder pressure and regulator setting in log book
\end{enumerate}

\subsubsection{Step 2: Power On SSM-5000A Main Unit}
\begin{enumerate}[noitemsep]
    \item Locate main power switch on rear panel (lower right corner)
    \item Switch to ON position (I)
    \item Observe front panel display:
    \begin{itemize}
        \item Display should illuminate within 5 seconds
        \item Initial message: "SHIMADZU SSM-5000A SYSTEM CHECK"
        \item Self-diagnostic runs for approximately 2--3 minutes
        \item Display will show "READY" when complete
    \end{itemize}
    \item Listen for internal fan operation (quiet humming sound)
    \item If error codes appear (E001, E002, etc.):
    \begin{itemize}
        \item Record error code
        \item Consult troubleshooting section (Section 11)
        \item Do not proceed until error is resolved
    \end{itemize}
\end{enumerate}

\subsubsection{Step 3: Start Computer and Software}
\begin{enumerate}[noitemsep]
    \item Power on dedicated analysis computer
    \item Wait for Windows operating system to fully load
    \item Locate SSM-5000A Control Software icon on desktop
    \item Double-click to launch software
    \item Software loading sequence:
    \begin{itemize}
        \item Splash screen appears (5 seconds)
        \item "Searching for instrument..." message displays
        \item Connection status window appears
        \item Green indicator shows successful connection
        \item If red indicator appears: check USB or RS-232 cable connection
    \end{itemize}
    \item Main software window opens with tabs: [Setup] [Calibration] [Sample] [Results]
    \item Verify software version matches instrument firmware (display on screen)
\end{enumerate}

\subsubsection{Step 4: Initialize Furnace System}
\begin{enumerate}[noitemsep]
    \item In software, navigate to [Setup] tab
    \item Click "Temperature Control" button
    \item Set TC furnace temperature:
    \begin{itemize}
        \item Click "TC Furnace" field
        \item Enter: \textbf{900°C}
        \item Click "Set" button
    \end{itemize}
    \item Set TIC furnace temperature:
    \begin{itemize}
        \item Click "TIC Furnace" field
        \item Enter: \textbf{200°C}
        \item Click "Set" button
    \end{itemize}
    \item Click "Start Heat-Up" button
    \item Monitor heating progress in status window:
    \begin{itemize}
        \item Current temperature displays in real-time
        \item TC furnace: heats at approximately 10°C/minute (90 minutes to 900°C)
        \item TIC furnace: heats at approximately 5°C/minute (40 minutes to 200°C)
    \end{itemize}
    \item \textbf{IMPORTANT:} Do not leave laboratory during heat-up
    \item Verify furnace status lights:
    \begin{itemize}
        \item Orange = heating
        \item Green = at set temperature
    \end{itemize}
\end{enumerate}

\subsubsection{Step 5: Initialize NDIR Detector}
\begin{enumerate}[noitemsep]
    \item Detector automatically powers on with main unit
    \item Navigate to [Setup] $>$ [Detector] in software
    \item Verify detector parameters:
    \begin{itemize}
        \item Detection range: 0--100\% CO$_2$ (auto-range)
        \item Response time: Standard (3 seconds)
        \item Zero adjustment: Auto
    \end{itemize}
    \item Allow 30 minutes minimum for detector stabilization
    \item Monitor baseline signal in real-time display window:
    \begin{itemize}
        \item Initial baseline may drift during warm-up
        \item Acceptable drift after stabilization: $<$ 5 ppm/hour
        \item Baseline noise: $<$ 2 ppm peak-to-peak
    \end{itemize}
    \item Perform automatic zero adjustment:
    \begin{itemize}
        \item Wait until temperature stability achieved (green indicator)
        \item Click "Auto Zero" button
        \item System purges detector cell with carrier gas (2 minutes)
        \item Baseline resets to zero
        \item Verify baseline stability within 5 minutes
    \end{itemize}
\end{enumerate}

\subsection{Gas Flow Rate Configuration}

\begin{enumerate}[noitemsep]
    \item Navigate to [Setup] $>$ [Gas Control] tab
    \item Configure TC analysis gas flows:
    \begin{itemize}
        \item Carrier gas (O$_2$): 500 mL/min
        \item Purge time before combustion: 30 seconds
        \item Combustion duration: 180 seconds (3 minutes)
    \end{itemize}
    \item Configure TIC analysis gas flows:
    \begin{itemize}
        \item Carrier gas (O$_2$): 300 mL/min
        \item Purge time before analysis: 20 seconds
        \item Analysis duration: 300 seconds (5 minutes)
    \end{itemize}
    \item Verify flow rates using:
    \begin{itemize}
        \item Software flow indicators (display in mL/min)
        \item OR external flow meter at gas outlet (optional verification)
    \end{itemize}
    \item Click "Apply Settings" to save configuration
\end{enumerate}

\subsection{System Blank Check}

\begin{enumerate}[noitemsep]
    \item Purpose: Verify system cleanliness and absence of contamination
    \item Pre-combust 5 ceramic boats:
    \begin{itemize}
        \item Place boats in TC furnace position
        \item Heat at 900°C for 10 minutes
        \item Cool in desiccator for 30 minutes minimum
    \end{itemize}
    \item Load 3 blank boats into autosampler positions 1--3
    \item In software, navigate to [Sample] tab
    \item Create blank sample sequence:
    \begin{itemize}
        \item Sample ID: BLANK-01, BLANK-02, BLANK-03
        \item Sample Type: Blank
        \item Weight: 0.0000 g (enter zero)
        \item Analysis: TC only
    \end{itemize}
    \item Click "Start Analysis" button
    \item System analyzes each blank boat automatically
    \item Review results in [Results] tab:
    \begin{itemize}
        \item Calculate average blank value
        \item Acceptable blank: $<$ 0.1 mg C per boat
        \item If blank $>$ 0.1 mg C:
        \begin{enumerate}
            \item Re-combust boats at 900°C for 15 minutes
            \item Retest blanks
            \item If still high: replace boats with new set
            \item Check for contamination in combustion tube
        \end{enumerate}
    \end{itemize}
    \item Record blank values in calibration log book
\end{enumerate}

\section{Standard Curve Creation and Verification}

\subsection{TC Calibration Curve Generation}

\subsubsection{Prepare Glucose Standard Solutions}
Follow Section 5.2 for preparation of 6 concentration levels.

\subsubsection{Load Standards into Software}
\begin{enumerate}[noitemsep]
    \item Navigate to [Calibration] tab in software
    \item Click "New Calibration" button
    \item Select calibration type: \textbf{TC (Total Carbon)}
    \item Click "Standard Table" button
    \item Enter standard information in table:
\end{enumerate}

\begin{table}[h]
\centering
\small
\begin{tabular}{cccccc}
\toprule
\textbf{Vial} & \textbf{Std ID} & \textbf{Conc.} & \textbf{Inj. Vol.} & \textbf{Expected} & \textbf{Replicates} \\
\textbf{Pos.} & & \textbf{(mg C/L)} & \textbf{($\mu$L)} & \textbf{Peak Area} & \\
\midrule
1 & TC-STD-1 & 10 & 100 & Low & 3 \\
2 & TC-STD-2 & 25 & 100 & Low-Med & 3 \\
3 & TC-STD-3 & 50 & 100 & Medium & 3 \\
4 & TC-STD-4 & 100 & 100 & Med-High & 3 \\
5 & TC-STD-5 & 250 & 100 & High & 3 \\
6 & TC-STD-6 & 500 & 100 & Very High & 3 \\
\bottomrule
\end{tabular}
\caption{TC Calibration Standard Entry Table}
\end{table}

\begin{enumerate}[noitemsep]
    \item For each standard level:
    \begin{itemize}
        \item Enter Standard ID in "Sample ID" field
        \item Enter concentration in "Concentration (mg C/L)" field
        \item Set injection volume to 100 $\mu$L
        \item Set number of replicates to 3
    \end{itemize}
    \item Click "Save Table" to store standard sequence
    \item Load standard vials into autosampler:
    \begin{itemize}
        \item Position standards in order (low to high concentration)
        \item Verify adequate volume in each vial ($>$ 500 $\mu$L)
        \item Check vial caps are properly sealed
        \item Ensure autosampler carousel is properly seated
    \end{itemize}
\end{enumerate}

\subsubsection{Run TC Calibration Sequence}
\begin{enumerate}[noitemsep]
    \item Verify furnaces at target temperature (green status indicators)
    \item Verify baseline stability ($<$ 5 ppm drift)
    \item Click "Start Calibration" button in software
    \item System performs automatic sequence:
    \begin{itemize}
        \item Injects 100 $\mu$L of standard solution onto pre-combusted boat
        \item Transfers boat to TC furnace (900°C)
        \item Combusts sample for 3 minutes in oxygen atmosphere
        \item CO$_2$ carried to NDIR detector
        \item Peak area integrated automatically
        \item Repeats for all 3 replicates
        \item Advances to next concentration level
    \end{itemize}
    \item Monitor analysis progress in real-time:
    \begin{itemize}
        \item Chromatogram displays each peak
        \item Peak area values appear in results table
        \item Verify peak shapes are symmetric (no shoulders or tailing)
        \item Verify retention times are consistent (within 5 seconds)
    \end{itemize}
    \item Total analysis time: approximately 90 minutes (6 levels $\times$ 3 reps $\times$ 5 min/sample)
\end{enumerate}

\subsubsection{Evaluate TC Calibration Curve}
\begin{enumerate}[noitemsep]
    \item After sequence completion, click "Calibration Results" button
    \item Software displays calibration curve graph:
    \begin{itemize}
        \item X-axis: Concentration (mg C/L)
        \item Y-axis: Peak Area (arbitrary units)
        \item Data points shown with error bars (standard deviation)
    \end{itemize}
    \item Review regression statistics displayed in results panel:
\end{enumerate}

\begin{table}[h]
\centering
\begin{tabular}{lcc}
\toprule
\textbf{Parameter} & \textbf{Specification} & \textbf{Action if Failed} \\
\midrule
Correlation coefficient ($r^2$) & $\geq$ 0.999 & Re-prepare standards \\
Residuals at each level & $<$ 10\% & Investigate outliers \\
Y-intercept & $<$ 5\% of lowest std response & Check blank contamination \\
Slope consistency & RSD $<$ 5\% over time & Verify detector stability \\
\bottomrule
\end{tabular}
\caption{TC Calibration Acceptance Criteria}
\end{table}

\begin{enumerate}[noitemsep]
    \item Calculate residuals for each standard:
    \begin{equation}
    \text{Residual (\%)} = \frac{|\text{Measured} - \text{Expected}|}{\text{Expected}} \times 100
    \end{equation}
    \item If any standard fails residual criterion:
    \begin{itemize}
        \item Mark outlier in software (right-click data point $>$ "Exclude")
        \item Re-analyze that concentration level only
        \item If repeat fails: remake that standard solution
    \end{itemize}
    \item If $r^2 < 0.999$:
    \begin{itemize}
        \item Check for non-linearity (may need quadratic fit)
        \item Verify standard preparation calculations
        \item Re-prepare entire standard series
        \item Check detector linearity with known reference standards
    \end{itemize}
    \item When calibration passes all criteria:
    \begin{itemize}
        \item Click "Accept Calibration" button
        \item Enter calibration ID: "TC-CAL-YYYYMMDD"
        \item Software saves calibration curve equation: $y = mx + b$
        \item Record slope (m), intercept (b), and $r^2$ in log book
        \item Print calibration report or save as PDF
    \end{itemize}
\end{enumerate}

\subsubsection{Calibration Documentation}
Create calibration file containing:
\begin{itemize}[noitemsep]
    \item Calibration curve graph with data points
    \item Regression equation and statistics
    \item Table of expected vs. measured values
    \item Peak chromatograms for all standards
    \item Analyst name, date, and instrument ID
    \item Lot numbers of standard materials used
\end{itemize}

\subsection{TIC Calibration Curve Generation}

\subsubsection{Prepare CaCO$_3$ Solid Standards}
\begin{enumerate}[noitemsep]
    \item Pre-combust ceramic boats at 900°C for 10 minutes
    \item Cool in desiccator for 30 minutes
    \item Tare analytical balance with empty boat
    \item Weigh CaCO$_3$ directly into boats according to Table 2
    \item Record exact weight to nearest 0.1 mg in log book
    \item Calculate expected carbon mass:
    \begin{equation}
    \text{C (mg)} = \text{CaCO}_3 \text{ weight (mg)} \times 0.120
    \end{equation}
    \item Prepare 3 boats at each concentration level (15 boats total)
    \item Load boats into autosampler carousel positions 1--15
\end{enumerate}

\subsubsection{Load TIC Standards into Software}
\begin{enumerate}[noitemsep]
    \item Navigate to [Calibration] tab
    \item Click "New Calibration" $>$ \textbf{TIC (Total Inorganic Carbon)}
    \item Click "Standard Table" button
    \item Enter standard information:
\end{enumerate}

\begin{table}[h]
\centering
\scriptsize
\begin{tabular}{ccccccc}
\toprule
\textbf{Vial} & \textbf{Std ID} & \textbf{CaCO$_3$} & \textbf{Carbon} & \textbf{Acid} & \textbf{Expected} & \textbf{Reps} \\
\textbf{Pos.} & & \textbf{(mg)} & \textbf{(mg)} & \textbf{($\mu$L)} & \textbf{Peak} & \\
\midrule
1--3 & TIC-STD-1 & 8.3 & 1.0 & 150 & Low & 3 \\
4--6 & TIC-STD-2 & 41.7 & 5.0 & 150 & Low-Med & 3 \\
7--9 & TIC-STD-3 & 83.3 & 10.0 & 200 & Medium & 3 \\
10--12 & TIC-STD-4 & 208.3 & 25.0 & 200 & Med-High & 3 \\
13--15 & TIC-STD-5 & 416.7 & 50.0 & 250 & High & 3 \\
\bottomrule
\end{tabular}
\caption{TIC Calibration Standard Entry Table}
\end{table}

\begin{enumerate}[noitemsep]
    \item For each replicate:
    \begin{itemize}
        \item Enter Sample ID (e.g., TIC-STD-1-A, TIC-STD-1-B, TIC-STD-1-C)
        \item Enter actual CaCO$_3$ weight measured
        \item Enter calculated carbon mass
        \item Specify acid volume to be added
    \end{itemize}
    \item Click "Save Table"
\end{enumerate}

\subsubsection{Run TIC Calibration Sequence}
\begin{enumerate}[noitemsep]
    \item Verify TIC furnace temperature: 200°C (green indicator)
    \item Verify acid reservoir contains fresh 25\% H$_3$PO$_4$ (minimum 50 mL)
    \item Prime acid delivery system:
    \begin{itemize}
        \item Click [Setup] $>$ [Acid System] $>$ "Prime"
        \item System draws acid through tubing (observe in waste bottle)
        \item Prime until no air bubbles visible in tubing (approximately 2 minutes)
        \item Click "Stop Prime"
    \end{itemize}
    \item Return to [Calibration] tab
    \item Click "Start Calibration" button
    \item System performs automatic sequence for each standard:
    \begin{enumerate}
        \item Autosampler positions boat at acid addition station
        \item Syringe dispenses specified volume of acid onto CaCO$_3$
        \item Reaction begins immediately: CaCO$_3$ + 2H$_3$PO$_4$ $\rightarrow$ 3CO$_2$ + Ca$_3$(PO$_4$)$_2$ + 3H$_2$O
        \item Wait 30 seconds for initial vigorous reaction
        \item Transfer boat to TIC furnace (200°C)
        \item Heat for 5 minutes to ensure complete carbonate decomposition
        \item Carrier gas sweeps CO$_2$ to NDIR detector
        \item Peak area integrated automatically
        \item Boat moved to waste position
        \item System advances to next standard
    \end{enumerate}
    \item Monitor analysis in real-time:
    \begin{itemize}
        \item Observe acid delivery (should be smooth, no sputtering)
        \item Watch for complete effervescence (bubbling stops within 1--2 minutes)
        \item Verify peak shapes are smooth and Gaussian
        \item Check for baseline return to zero between samples
    \end{itemize}
    \item Total analysis time: approximately 120 minutes (15 samples $\times$ 8 min/sample)
\end{enumerate}

\subsubsection{Evaluate TIC Calibration Curve}
\begin{enumerate}[noitemsep]
    \item Click "Calibration Results" after sequence completion
    \item Review calibration curve graph (mg C vs. Peak Area)
    \item Examine regression statistics:
\end{enumerate}

\begin{table}[h]
\centering
\begin{tabular}{lcc}
\toprule
\textbf{Parameter} & \textbf{Specification} & \textbf{Action if Failed} \\
\midrule
Correlation coefficient ($r^2$) & $\geq$ 0.998 & Re-run calibration \\
Residuals at each level & $<$ 15\% & Check acid addition \\
Y-intercept & $<$ 10\% of lowest std & Check system blank \\
Replicate precision (RSD) & $<$ 10\% & Improve weighing precision \\
Recovery at mid-level & 90--110\% & Verify CaCO$_3$ purity \\
\bottomrule
\end{tabular}
\caption{TIC Calibration Acceptance Criteria}
\end{table}

\begin{enumerate}[noitemsep]
    \item Common issues and troubleshooting:
    \begin{itemize}
        \item \textbf{High variability at low levels:} Increase sample size or acid volume
        \item \textbf{Incomplete reaction:} Extend heating time to 7--10 minutes
        \item \textbf{Poor peak shape:} Check for acid delivery blockage, clean syringe
        \item \textbf{Low recovery:} Verify CaCO$_3$ dried properly, check for moisture uptake
    \end{itemize}
    \item When calibration passes all criteria:
    \begin{itemize}
        \item Click "Accept Calibration"
        \item Enter calibration ID: "TIC-CAL-YYYYMMDD"
        \item Save calibration file with all documentation
        \item Print calibration report for laboratory records
    \end{itemize}
\end{enumerate}

\subsection{Calibration Verification with QC Standards}

\subsubsection{Prepare QC Check Standards}
\begin{enumerate}[noitemsep]
    \item \textbf{TC QC Standard:} Prepare mid-level glucose standard (100 mg C/L)
    \begin{itemize}
        \item Use different glucose lot than calibration standards (if available)
        \item Prepare fresh from stock solution
        \item Target: expected recovery 95--105\%
    \end{itemize}
    \item \textbf{TIC QC Standard:} Weigh 83.3 mg CaCO$_3$ (10 mg C equivalent)
    \begin{itemize}
        \item Use different CaCO$_3$ lot than calibration (if available)
        \item Prepare 3 replicates
        \item Target: expected recovery 90--110\%
    \end{itemize}
    \item \textbf{Certified Reference Material (CRM):} If available
    \begin{itemize}
        \item Use certified soil standard with known TC and TIC
        \item Examples: NIST SRM 2709a (San Joaquin Soil), NIST SRM 2711a (Montana Soil)
        \item Prepare 3 replicates at recommended weight
    \end{itemize}
\end{enumerate}

\subsubsection{Analyze QC Standards}
\begin{enumerate}[noitemsep]
    \item Navigate to [Sample] tab
    \item Create QC sample sequence:
\end{enumerate}

\begin{table}[h]
\centering
\begin{tabular}{lcccc}
\toprule
\textbf{Sample ID} & \textbf{Type} & \textbf{Expected} & \textbf{Accept. Range} & \textbf{Replicates} \\
\midrule
TC-QC-MID & TC Check & 100 mg C/L & 95--105 mg C/L & 3 \\
TIC-QC-MID & TIC Check & 10.0 mg C & 9.0--11.0 mg C & 3 \\
CRM-2709a & Reference Soil & Per certificate & Per certificate & 3 \\
\bottomrule
\end{tabular}
\caption{QC Standard Verification Sequence}
\end{table}

\begin{enumerate}[noitemsep]
    \item Load QC samples into autosampler
    \item Click "Start Analysis"
    \item System analyzes using active calibration curves
    \item Review results and calculate recovery:
    \begin{equation}
    \text{Recovery (\%)} = \frac{\text{Measured Value}}{\text{Expected Value}} \times 100
    \end{equation}
    \item QC acceptance criteria:
    \begin{itemize}
        \item TC QC recovery: 95--105\%
        \item TIC QC recovery: 90--110\%
        \item CRM recovery: within certified uncertainty range
        \item Replicate RSD: $<$ 10\%
    \end{itemize}
    \item If QC fails:
    \begin{itemize}
        \item Re-analyze QC to rule out injection/analysis error
        \item If repeat fails: investigate calibration
        \item Possible causes: standard preparation error, detector drift, contamination
        \item Action: Recalibrate system and re-analyze QC
        \item Do NOT proceed to sample analysis until QC passes
    \end{itemize}
    \item Record QC results in Quality Control Log:
    \begin{itemize}
        \item Date and time
        \item QC ID and lot number
        \item Expected and measured values
        \item Recovery percentage
        \item Analyst initials
        \item Pass/Fail designation
    \end{itemize}
\end{enumerate}

\section{Sample Table Preparation and Analysis}

\subsection{Create Sample Analysis Sequence}

\subsubsection{Sample Organization Strategy}
\begin{enumerate}[noitemsep]
    \item Organize samples by expected carbonate content (high, medium, low)
    \item Analyze samples with similar carbonate levels together
    \item Include appropriate QC samples throughout sequence
    \item Recommended sequence structure:
\end{enumerate}

\begin{longtable}{cccccc}
\toprule
\textbf{Position} & \textbf{Sample ID} & \textbf{Type} & \textbf{Analysis} & \textbf{Weight} & \textbf{Reps} \\
\midrule
\endfirsthead
\multicolumn{6}{c}%
{\tablename\ \thetable\ -- \textit{Continued from previous page}} \\
\toprule
\textbf{Position} & \textbf{Sample ID} & \textbf{Type} & \textbf{Analysis} & \textbf{Weight} & \textbf{Reps} \\
\midrule
\endhead
\midrule
\multicolumn{6}{r}{\textit{Continued on next page}} \\
\endfoot
\bottomrule
\endlastfoot
1 & BLANK-01 & Method Blank & TC & 0.0000 g & 1 \\
2--4 & TC-QC-MID & QC Standard & TC & 100 mg/L & 3 \\
5--7 & SOIL-001-A & Unknown & TC & 50.0 mg & 3 \\
8--10 & SOIL-001-B & Unknown & TC & 50.0 mg & 3 \\
11--13 & SOIL-002-A & Unknown & TC & 45.2 mg & 3 \\
14--16 & SOIL-002-B & Unknown & TC & 45.2 mg & 3 \\
17--19 & SOIL-003-A & Unknown & TC & 52.8 mg & 3 \\
20--22 & SOIL-003-B & Unknown & TC & 52.8 mg & 3 \\
23 & BLANK-02 & Method Blank & TC & 0.0000 g & 1 \\
24--26 & TC-QC-MID & QC Standard & TC & 100 mg/L & 3 \\
27--29 & SOIL-004-A & Unknown & TC & 38.5 mg & 3 \\
\multicolumn{6}{c}{... continue sample sequence ...} \\
\end{longtable}

\subsubsection{Enter Sample Information into Software}
\begin{enumerate}[noitemsep]
    \item Navigate to [Sample] tab in software
    \item Click "New Sequence" button
    \item Enter sequence information:
    \begin{itemize}
        \item Sequence Name: "PROJECT-YYYYMMDD-TC"
        \item Analyst Name: [Your name]
        \item Project ID: [If applicable]
        \item Analysis Type: TC or TIC
    \end{itemize}
    \item Click "Sample Table" button
    \item For each sample position, enter:
    \begin{itemize}
        \item Vial Position: 1--40 (autosampler position)
        \item Sample ID: Unique identifier
        \item Sample Type: Unknown/QC/Blank/Duplicate
        \item Sample Weight: Exact weight in mg (dry weight basis)
        \item Dilution Factor: Usually 1 (adjust if sample pre-diluted)
        \item Number of Replicates: 3 (or 5 for low-level samples)
        \item Comments: Sample description, carbonate level, special notes
    \end{itemize}
    \item Use software features:
    \begin{itemize}
        \item Copy/Paste function for duplicate entries
        \item "Fill Down" to replicate information
        \item Import from Excel spreadsheet (if available)
    \end{itemize}
    \item Review entire table for errors before saving
    \item Click "Save Sequence" 
    \item File naming convention: "PROJ-ANALYST-YYYYMMDD-TCOR TIC.seq"
\end{enumerate}

\subsection{Sample Weighing and Loading}

\subsubsection{Prepare Sample Boats}
\begin{enumerate}[noitemsep]
    \item Pre-combust ceramic boats:
    \begin{itemize}
        \item Load 40 boats (full carousel) into TC furnace
        \item Heat at 900°C for 10 minutes
        \item Transfer immediately to desiccator using heat-resistant tongs
        \item Cool for minimum 30 minutes
    \end{itemize}
    \item Label boats with sample IDs using:
    \begin{itemize}
        \item High-temperature ceramic marker, OR
        \item Small adhesive labels (note: will burn off during analysis)
        \item Match boat labels to autosampler positions
    \end{itemize}
\end{enumerate}

\subsubsection{Weigh Soil Samples}
\begin{enumerate}[noitemsep]
    \item Set up analytical balance in draft-free area
    \item Place anti-static mat on balance pan
    \item Calibrate balance with certified weights if not auto-calibrating
    \item For each sample:
    \begin{enumerate}
        \item Close balance doors
        \item Press "Tare" button with empty boat on pan
        \item Verify zero reading (0.0000 g)
        \item Open door slowly to prevent air currents
        \item Add soil sample using micro-spatula:
        \begin{itemize}
            \item Add sample in small increments
            \item Target weight according to expected carbonate content
            \item High carbonate: 10--20 mg
            \item Medium carbonate: 20--40 mg
            \item Low carbonate: 40--80 mg
        \end{itemize}
        \item Close balance door and wait for stable reading
        \item Record weight when stability indicator appears ($<$ 5 seconds)
        \item Record weight to 0.0001 g precision in laboratory notebook
        \item Transfer boat immediately to labeled position in autosampler
    \end{enumerate}
    \item Quality checks during weighing:
    \begin{itemize}
        \item Verify boat placement in center of pan
        \item Check for powder on boat exterior (clean if necessary)
        \item Ensure sample is evenly distributed in boat
        \item Avoid touching boat with bare hands (use powder-free gloves or tweezers)
    \end{itemize}
    \item Cross-reference:
    \begin{itemize}
        \item Check sample ID on boat matches sequence table
        \item Verify autosampler position corresponds to sequence
        \item Confirm weight is within acceptable range for carbonate level
    \end{itemize}
\end{enumerate}

\subsubsection{Load Autosampler Carousel}
\begin{enumerate}[noitemsep]
    \item Remove carousel from instrument:
    \begin{itemize}
        \item Open autosampler door
        \item Release carousel lock mechanism
        \item Lift carousel straight up and out
        \item Place on clean, level surface
    \end{itemize}
    \item Load boats into numbered positions:
    \begin{itemize}
        \item Match boat to position number in sequence table
        \item Ensure boats are fully seated in holders
        \item Check boat clips engage properly
        \item Verify no boats are tilted or unstable
    \end{itemize}
    \item Visual inspection:
    \begin{itemize}
        \item Count total boats loaded (should match sequence length)
        \item Verify empty positions correspond to sequence gaps
        \item Check for any damaged or cracked boats
    \end{itemize}
    \item Reinstall carousel:
    \begin{itemize}
        \item Align carousel orientation marks
        \item Lower carousel into instrument
        \item Engage lock mechanism (hear/feel click)
        \item Manually rotate carousel to verify smooth operation
        \item Close autosampler door
    \end{itemize}
    \item Software verification:
    \begin{itemize}
        \item In software, click "Check Carousel"
        \item System performs position check (rotates to each position)
        \item Verify "Carousel OK" message appears
        \item If error: reload carousel and recheck
    \end{itemize}
\end{enumerate}

\subsection{TC Analysis Execution}

\subsubsection{Pre-Analysis Checklist}
Verify all items before starting:
\begin{itemize}[noitemsep]
    \item[$\Box$] TC furnace at 900°C (green indicator)
    \item[$\Box$] Baseline stable ($<$ 5 ppm drift)
    \item[$\Box$] Oxygen pressure 0.45 MPa (65 psi)
    \item[$\Box$] Calibration current and verified
    \item[$\Box$] QC standards passed acceptance criteria
    \item[$\Box$] Sample sequence saved
    \item[$\Box$] All samples weighed and loaded correctly
    \item[$\Box$] Carousel installed and operational
    \item[$\Box$] Waste containers not full
    \item[$\Box$] Computer hard drive has adequate space ($>$ 1 GB)
\end{itemize}

\subsubsection{Start Analysis Sequence}
\begin{enumerate}[noitemsep]
    \item In [Sample] tab, select saved sequence file
    \item Review sequence summary:
    \begin{itemize}
        \item Total number of samples
        \item Estimated run time
        \item Calibration curve to be used
    \end{itemize}
    \item Click "Start Analysis" button
    \item Confirmation dialog appears:
    \begin{itemize}
        \item Verify sequence name correct
        \item Check starting position (usually Position 1)
        \item Confirm analysis parameters
        \item Click "OK" to proceed
    \end{itemize}
    \item System begins automated analysis:
    \begin{enumerate}
        \item Carousel rotates to first sample position
        \item Boat transferred to loading zone
        \item Boat moved into TC furnace (900°C)
        \item Oxygen flow purges system (30 seconds)
        \item Combustion period begins (180 seconds):
        \begin{itemize}
            \item All carbon oxidized to CO$_2$
            \item CO$_2$ swept to NDIR detector by carrier gas
            \item Real-time chromatogram displays on screen
        \end{itemize}
        \item Peak detection and integration:
        \begin{itemize}
            \item Software identifies peak start/end points
            \item Integrates peak area
            \item Calculates carbon concentration using calibration curve
        \end{itemize}
        \item Post-combustion purge (30 seconds)
        \item Boat moved to cooling zone
        \item System advances to next sample automatically
    \end{enumerate}
\end{enumerate}

\subsubsection{Monitor Analysis Progress}
\begin{enumerate}[noitemsep]
    \item Real-time monitoring displays:
    \begin{itemize}
        \item Current sample ID and position
        \item Live chromatogram with peak development
        \item Current sample number / total samples
        \item Estimated time remaining
        \item Last 5 results with statistics
    \end{itemize}
    \item Quality indicators to watch:
    \begin{itemize}
        \item Peak shapes: should be smooth, Gaussian, symmetric
        \item Retention times: should be consistent ($\pm$ 5 seconds)
        \item Baseline stability: should return to zero between peaks
        \item Replicate agreement: RSD should be $<$ 10\%
    \end{itemize}
    \item Alert conditions:
    \begin{itemize}
        \item Software displays warning messages for:
        \begin{enumerate}
            \item Peak area outside calibration range
            \item Poor peak shape (excessive tailing)
            \item High RSD between replicates ($>$ 15\%)
            \item Baseline drift exceeds limits
        \end{enumerate}
        \item Investigate warnings immediately
        \item Pause sequence if necessary to resolve issues
    \end{itemize}
    \item \textbf{Important:} Do not leave instrument unattended
    \begin{itemize}
        \item Check progress every 30--60 minutes
        \item Be present at laboratory during entire run if possible
        \item If must leave: ensure another trained analyst monitors
    \end{itemize}
\end{enumerate}

\subsection{TIC Analysis Execution}

After completing TC analysis, perform TIC analysis on same samples:

\subsubsection{Prepare TIC Sample Boats}
\begin{enumerate}[noitemsep]
    \item Pre-combust fresh boats at 900°C (as before)
    \item Cool in desiccator for 30 minutes
    \item Weigh same soil samples into new boats:
    \begin{itemize}
        \item Use identical sample masses as TC analysis
        \item Record weights to 0.0001 g
        \item Label boats with sample IDs + "-TIC" designation
    \end{itemize}
\end{enumerate}

\subsubsection{Acid Addition}
\begin{enumerate}[noitemsep]
    \item Verify 25\% H$_3$PO$_4$ in acid reservoir ($>$ 50 mL)
    \item Prime acid delivery system (if first use of day)
    \item Load boats into autosampler carousel
    \item Create TIC sample sequence in software:
    \begin{itemize}
        \item Use same sample IDs with "-TIC" suffix
        \item Enter same sample weights
        \item Specify acid volume for each sample:
        \begin{enumerate}
            \item High carbonate (vigorous reaction): 150--200 $\mu$L
            \item Medium carbonate: 150 $\mu$L
            \item Low carbonate: 100--150 $\mu$L
        \end{enumerate}
    \end{itemize}
\end{enumerate}

\subsubsection{Start TIC Analysis}
\begin{enumerate}[noitemsep]
    \item Verify TIC furnace at 200°C (green indicator)
    \item Select TIC calibration curve
    \item Click "Start Analysis"
    \item System performs for each sample:
    \begin{enumerate}
        \item Positions boat at acid addition station
        \item Dispenses specified volume of H$_3$PO$_4$
        \item Waits 30 seconds for initial reaction
        \item Transfers boat to TIC furnace (200°C)
        \item Heats for 5--7 minutes
        \item \textbf{Reaction:} Carbonates decompose to CO$_2$ at low temperature
        \item Organic carbon remains unoxidized (requires 900°C)
        \item CO$_2$ swept to detector and quantified
    \end{enumerate}
    \item Monitor for complete carbonate decomposition:
    \begin{itemize}
        \item Peak should reach plateau and return to baseline
        \item Incomplete reaction shows gradual tail
        \item If incomplete: increase heating time to 10 minutes
    \end{itemize}
\end{enumerate}

\section{Standard Curve Verification and Quality Control}

\subsection{Within-Run QC Checks}

\subsubsection{Continuing Calibration Verification (CCV)}
\begin{enumerate}[noitemsep]
    \item Analyze QC check standard every 10 samples:
    \begin{itemize}
        \item Insert QC standard into sequence at regular intervals
        \item Use mid-level standard (e.g., 100 mg C/L for TC, 10 mg C for TIC)
        \item Prepare QC from separate stock solution (independent verification)
    \end{itemize}
    \item Acceptance criteria for CCV:
    \begin{itemize}
        \item Recovery: 90--110\% of expected value
        \item RSD of triplicates: $<$ 10\%
        \item If QC fails:
        \begin{enumerate}
            \item Flag affected samples in database
            \item Re-analyze QC to confirm failure
            \item If repeat fails: stop sequence
            \item Investigate cause: detector drift, contamination, standard degradation
            \item Recalibrate if necessary
            \item Re-analyze all samples since last passing QC
        \end{enumerate}
    \end{itemize}
    \item Document each QC result:
    \begin{itemize}
        \item Record in QC Control Chart
        \item Plot recovery vs. time
        \item Look for trends (systematic drift)
        \item Calculate running statistics (mean, SD, control limits)
    \end{itemize}
\end{enumerate}

\subsubsection{Method Blanks}
\begin{enumerate}[noitemsep]
    \item Analyze blank boat every 10 samples:
    \begin{itemize}
        \item Pre-combusted ceramic boat with no sample
        \item Analyze using same method as samples
        \item Expected result: $<$ 0.1 mg C
    \end{itemize}
    \item Blank acceptance:
    \begin{itemize}
        \item Blank $<$ 10\% of lowest sample concentration: acceptable
        \item Blank $>$ 10\% of sample: investigate contamination
        \item Possible causes: dirty boats, contaminated gases, carryover
    \end{itemize}
    \item Corrective actions if blank fails:
    \begin{itemize}
        \item Re-combust boats at 900°C for 15 minutes
        \item Replace quartz wool in combustion tube
        \item Check for leaks in gas lines
        \item Verify gas purity specifications
    \end{itemize}
\end{enumerate}

\subsubsection{Duplicate Samples}
\begin{enumerate}[noitemsep]
    \item Analyze 10\% of samples as duplicates:
    \begin{itemize}
        \item Prepare duplicate from same homogenized soil
        \item Weigh independently
        \item Analyze as separate sample
        \item Label as "-DUP" in sample ID
    \end{itemize}
    \item Calculate Relative Percent Difference (RPD):
    \begin{equation}
    \text{RPD} = \frac{|\text{Sample} - \text{Duplicate}|}{\left(\frac{\text{Sample} + \text{Duplicate}}{2}\right)} \times 100
    \end{equation}
    \item Acceptance criteria:
    \begin{itemize}
        \item RPD $<$ 15\% for samples $>$ 0.5\% C
        \item RPD $<$ 20\% for samples $<$ 0.5\% C
        \item If RPD exceeds limit:
        \begin{enumerate}
            \item Check sample homogeneity
            \item Re-grind and mix sample
            \item Analyze third replicate
            \item If variability persists: report as heterogeneous sample
        \end{enumerate}
    \end{itemize}
\end{enumerate}

\subsection{Calibration Curve Stability Assessment}

\subsubsection{Daily Calibration Check}
\begin{enumerate}[noitemsep]
    \item At start of each analysis day:
    \begin{itemize}
        \item Analyze 3 calibration levels (low, mid, high)
        \item Calculate recovery for each level
        \item Compare to original calibration curve
    \end{itemize}
    \item Acceptable drift:
    \begin{itemize}
        \item All levels recover within 95--105\%
        \item Correlation coefficient remains $\geq$ 0.999 for TC, $\geq$ 0.998 for TIC
        \item Slope change $<$ 10\% from original calibration
    \end{itemize}
    \item If drift exceeds limits:
    \begin{itemize}
        \item Perform full recalibration with all standards
        \item Investigate cause: detector aging, gas flow changes, contamination
        \item Service instrument if necessary
    \end{itemize}
\end{enumerate}

\subsubsection{Calibration Frequency Requirements}
\begin{itemize}[noitemsep]
    \item Full calibration (6 levels minimum): Every 7 days or 100 samples
    \item Calibration verification: Beginning of each analysis day
    \item CCV check standards: Every 10 samples during run
    \item Recalibration required if:
    \begin{itemize}
        \item QC recovery outside 90--110\%
        \item Detector or furnace maintenance performed
        \item Gas cylinder replaced
        \item More than 14 days since last calibration
    \end{itemize}
\end{itemize}

\subsubsection{Calibration Curve Documentation}
Maintain calibration file containing:
\begin{enumerate}[noitemsep]
    \item Calibration date and analyst name
    \item Standard preparation records:
    \begin{itemize}
        \item Chemical lot numbers
        \item Weighing records
        \item Dilution calculations
        \item Expiration dates
    \end{itemize}
    \item Raw data:
    \begin{itemize}
        \item Individual replicate peak areas
        \item Average peak areas per level
        \item Standard deviations and RSD
    \end{itemize}
    \item Calibration curve graph (print or save as PDF)
    \item Regression statistics:
    \begin{itemize}
        \item Equation: $y = mx + b$
        \item Slope (m), intercept (b)
        \item Correlation coefficient ($r^2$)
        \item Residuals at each level
    \end{itemize}
    \item QC verification results
    \item Acceptance/rejection decision with signature
    \item File in laboratory records for minimum 5 years
\end{enumerate}

\section{Data Recording and Calculations}

\subsection{Raw Data Collection}

Record in bound laboratory notebook or electronic laboratory notebook (ELN):

\subsubsection{Sample Information}
\begin{itemize}[noitemsep]
    \item Sample ID and description
    \item Sample collection date and location
    \item Sample preparation date
    \item Moisture content (with calculation)
    \item Sample weight (as-received and dry basis)
    \item Date and time of analysis
\end{itemize}

\subsubsection{Instrument Information}
\begin{itemize}[noitemsep]
    \item Instrument ID: Shimadzu SSM-5000A Serial Number
    \item Analyst name and initials
    \item Calibration ID and date
    \item Calibration verification results
    \item Gas cylinder lot numbers and pressures
    \item Reagent lot numbers (phosphoric acid, standards)
\end{itemize}

\subsubsection{Analysis Data}
\begin{itemize}[noitemsep]
    \item Raw peak areas for each replicate
    \item Calculated carbon concentrations (mg C)
    \item Mean, standard deviation, RSD
    \item QC check results
    \item Blank values
    \item Any deviations from SOP
\end{itemize}

\subsection{Detailed Calculation Procedures}

\subsubsection{Moisture Content Calculation}
\begin{equation}
\text{Moisture (\%)} = \frac{\text{Wet Weight}_{\text{soil}} - \text{Dry Weight}_{\text{soil}}}{\text{Wet Weight}_{\text{soil}}} \times 100
\end{equation}

\textbf{Example:}
\begin{itemize}[noitemsep]
    \item Crucible weight: 25.3421 g
    \item Crucible + wet soil: 28.5632 g
    \item Wet soil weight: 28.5632 $-$ 25.3421 = 3.2211 g
    \item Crucible + dry soil (after 105°C): 28.2156 g
    \item Dry soil weight: 28.2156 $-$ 25.3421 = 2.8735 g
    \item Moisture: $\frac{3.2211 - 2.8735}{3.2211} \times 100 = 10.79\%$
\end{itemize}

\subsubsection{Dry Weight Correction}
\begin{equation}
\text{Dry Weight}_{\text{soil}} = \text{Wet Weight}_{\text{soil}} \times \left(1 - \frac{\text{Moisture \%}}{100}\right)
\end{equation}

\textbf{Example:}
\begin{itemize}[noitemsep]
    \item Sample weighed for analysis: 52.3 mg (as-received moisture content)
    \item Moisture content: 10.79\%
    \item Dry weight: $52.3 \times (1 - 0.1079) = 52.3 \times 0.8921 = 46.7$ mg
\end{itemize}

\subsubsection{Carbon Concentration from Calibration}

Software calculates automatically using calibration equation, but manual check:

\begin{equation}
\text{Carbon Mass (mg)} = \frac{\text{Peak Area} - b}{m}
\end{equation}

where $m$ = slope, $b$ = intercept from calibration curve

\textbf{Example TC Analysis:}
\begin{itemize}[noitemsep]
    \item Peak area: 125,432 (arbitrary units)
    \item Calibration slope ($m$): 2,508.6 AU per mg C
    \item Calibration intercept ($b$): 1,245 AU
    \item Carbon mass: $\frac{125,432 - 1,245}{2,508.6} = 49.5$ mg C
\end{itemize}

\subsubsection{Carbon Percentage Calculation}
\begin{equation}
\text{Carbon (\%)} = \frac{\text{Carbon Mass (mg)}}{\text{Sample Dry Weight (mg)}} \times 100
\end{equation}

\textbf{Example:}
\begin{itemize}[noitemsep]
    \item TC measured: 49.5 mg C
    \item Sample dry weight: 46.7 mg
    \item TC (\%): $\frac{49.5}{46.7} \times 100 = 106.0\%$ 
    \item \textbf{Note:} This result $>$100\% indicates error - investigate!
    \item Likely cause: incorrect sample weight or moisture calculation
\end{itemize}

\textbf{Corrected Example:}
\begin{itemize}[noitemsep]
    \item TC measured: 2.35 mg C
    \item Sample dry weight: 46.7 mg
    \item TC (\%): $\frac{2.35}{46.7} \times 100 = 5.03\%$
\end{itemize}

\subsubsection{TOC Calculation}
\begin{equation}
\text{TOC (\%)} = \text{TC (\%)} - \text{TIC (\%)}
\end{equation}

\textbf{Example:}
\begin{itemize}[noitemsep]
    \item TC: 5.03\%
    \item TIC: 2.18\%
    \item TOC: $5.03 - 2.18 = 2.85\%$
\end{itemize}

\textbf{Special Case - Negative TOC:}
\begin{itemize}[noitemsep]
    \item If TIC $>$ TC by small amount ($<$ 0.5\%): measurement uncertainty
    \item Report TOC as "Below Detection Limit" or "0.0\%"
    \item If TIC $>>$ TC: serious error in measurement - reanalyze
\end{itemize}

\subsubsection{Statistical Calculations}

\textbf{Mean of Replicates:}
\begin{equation}
\overline{x} = \frac{\sum_{i=1}^{n} x_i}{n}
\end{equation}

\textbf{Standard Deviation:}
\begin{equation}
s = \sqrt{\frac{\sum_{i=1}^{n}(x_i - \overline{x})^2}{n-1}}
\end{equation}

\textbf{Relative Standard Deviation (RSD):}
\begin{equation}
\text{RSD (\%)} = \frac{s}{\overline{x}} \times 100
\end{equation}

\textbf{Example Calculation:}

Three TC replicates: 5.12\%, 4.98\%, 5.07\%

\begin{itemize}[noitemsep]
    \item Mean: $\overline{x} = \frac{5.12 + 4.98 + 5.07}{3} = 5.06\%$
    \item Deviations: 0.06, $-$0.08, 0.01
    \item Variance: $\frac{(0.06)^2 + (-0.08)^2 + (0.01)^2}{3-1} = \frac{0.0101}{2} = 0.00505$
    \item Standard deviation: $s = \sqrt{0.00505} = 0.071\%$
    \item RSD: $\frac{0.071}{5.06} \times 100 = 1.4\%$ (acceptable, $<$10\%)
\end{itemize}

\subsubsection{Relative Percent Difference (RPD)}
For duplicate samples:
\begin{equation}
\text{RPD (\%)} = \frac{|x_1 - x_2|}{\left(\frac{x_1 + x_2}{2}\right)} \times 100
\end{equation}

\textbf{Example:}
\begin{itemize}[noitemsep]
    \item Sample: 5.06\% TC
    \item Duplicate: 4.89\% TC
    \item RPD: $\frac{|5.06 - 4.89|}{\frac{5.06 + 4.89}{2}} \times 100 = \frac{0.17}{4.975} \times 100 = 3.4\%$ (acceptable)
\end{itemize}

\subsection{Significant Figures and Rounding}

\begin{itemize}[noitemsep]
    \item Sample weights: Record to 0.0001 g (4 decimal places)
    \item Peak areas: Software reports to nearest integer
    \item Carbon concentrations: Report to 2 decimal places
    \item Final TC, TIC, TOC results: Report to 2 decimal places
    \item Standard deviations: Report to same decimal places as mean
    \item Percent values: Report to 0.01\% (2 decimal places)
\end{itemize}

\textbf{Rounding Rules:}
\begin{itemize}[noitemsep]
    \item If digit to right is $<$ 5: round down
    \item If digit to right is $>$ 5: round up
    \item If digit to right is exactly 5: round to nearest even number
    \item Carry full precision through calculations, round only final result
\end{itemize}

\subsection{Reporting Results}

\subsubsection{Standard Report Format}

For each sample, report:

\begin{table}[h]
\centering
\small
\begin{tabular}{lcccc}
\toprule
\textbf{Sample ID} & \textbf{TC (\%)} & \textbf{TIC (\%)} & \textbf{TOC (\%)} & \textbf{Moisture (\%)} \\
\midrule
SOIL-001 & 5.06 $\pm$ 0.07 & 2.18 $\pm$ 0.12 & 2.88 $\pm$ 0.14 & 10.8 \\
SOIL-002 & 3.42 $\pm$ 0.05 & 0.85 $\pm$ 0.08 & 2.57 $\pm$ 0.09 & 8.3 \\
SOIL-003 & 1.95 $\pm$ 0.03 & 0.12 $\pm$ 0.02 & 1.83 $\pm$ 0.04 & 5.2 \\
\bottomrule
\end{tabular}
\caption{Example Results Report (mean $\pm$ standard deviation, $n=3$, dry weight basis)}
\end{table}

\subsubsection{Data Qualifiers}

Use appropriate data qualifiers when necessary:

\begin{itemize}[noitemsep]
    \item \textbf{J} = Estimated value (QC slightly outside limits but usable)
    \item \textbf{U} = Below detection limit
    \item \textbf{D} = Diluted sample
    \item \textbf{*} = Duplicate RPD exceeded but values reported
    \item \textbf{ND} = Not detected
    \item \textbf{NA} = Not analyzed (e.g., TIC not performed on non-calcareous soil)
\end{itemize}

\subsubsection{Required Documentation}

Include with every report:
\begin{itemize}[noitemsep]
    \item Sample identification and description
    \item Analysis date(s)
    \item Method reference (this SOP number)
    \item Analyst name
    \item Instrument ID
    \item Calibration information
    \item QC results summary
    \item Moisture content
    \item Reporting basis (dry weight)
    \item Any deviations or qualifiers
    \item Uncertainty estimates or confidence intervals
\end{itemize}

\section{Maintenance and Troubleshooting}

\subsection{Daily Maintenance}

Perform at end of each analysis day:

\begin{enumerate}[noitemsep]
    \item \textbf{Check gas pressures and flows:}
    \begin{itemize}
        \item Record oxygen cylinder pressure in log
        \item Verify regulator maintains 0.45 MPa
        \item Check for pressure drops indicating leaks
    \end{itemize}
    \item \textbf{Verify furnace temperatures:}
    \begin{itemize}
        \item TC furnace: 900°C $\pm$ 5°C
        \item TIC furnace: 200°C $\pm$ 5°C
        \item Record actual temperatures in maintenance log
    \end{itemize}
    \item \textbf{Inspect sample inlet and combustion tube:}
    \begin{itemize}
        \item Look for sample residue buildup
        \item Check quartz wool color (replace if brown/black)
        \item Verify no cracks or damage to combustion tube
    \end{itemize}
    \item \textbf{Clean ceramic boats:}
    \begin{itemize}
        \item Remove used boats from waste position
        \item Inspect for damage (cracks, chips)
        \item Discard damaged boats
        \item Pre-combust reusable boats at 900°C (10 min)
        \item Store cleaned boats in desiccator
    \end{itemize}
    \item \textbf{Empty waste containers:}
    \begin{itemize}
        \item Check acid waste bottle level
        \item Empty if $>$ 50\% full
        \item Check condensate trap
    \end{itemize}
    \item \textbf{Check desiccant:}
    \begin{itemize}
        \item Observe indicating desiccant color
        \item Blue/orange = active (acceptable)
        \item Pink/white = saturated (replace immediately)
    \end{itemize}
    \item \textbf{Backup data:}
    \begin{itemize}
        \item Save all analysis files to network drive
        \item Export results to Excel or LIMS
        \item Create backup copy on external drive
    \end{itemize}
\end{enumerate}

\subsection{Weekly Maintenance}

\begin{enumerate}[noitemsep]
    \item \textbf{Clean combustion tubes thoroughly:}
    \begin{itemize}
        \item Allow furnaces to cool to $<$ 100°C
        \item Remove combustion tubes carefully (wear heat-resistant gloves)
        \item Inspect for carbon deposits or white residue (carbonates)
        \item Clean with appropriate brush and deionized water
        \item Dry completely in oven at 105°C
        \item Replace quartz wool packing (fresh, pre-combusted)
        \item Reinstall tubes with new O-rings if needed
    \end{itemize}
    \item \textbf{Clean NDIR detector cell:}
    \begin{itemize}
        \item Follow manufacturer procedure for cell cleaning
        \item Use methanol and lint-free wipes
        \item Check detector windows for fogging or deposits
        \item Allow to dry completely before reassembly
        \item Perform detector zero adjustment after cleaning
    \end{itemize}
    \item \textbf{Verify calibration stability:}
    \begin{itemize}
        \item Analyze 3 calibration levels
        \item Compare to stored calibration curve
        \item Document any drift
        \item Recalibrate if needed
    \end{itemize}
    \item \textbf{Check acid delivery system:}
    \begin{itemize}
        \item Inspect syringe for bubbles or blockage
        \item Clean syringe barrel and needle
        \item Prime system thoroughly
        \item Test delivery volume accuracy (dispense onto weighing paper, weigh)
        \item Adjust volume calibration if necessary
    \end{itemize}
    \item \textbf{Leak test gas lines:}
    \begin{itemize}
        \item Apply leak detection solution to all fittings
        \item Pressurize system
        \item Look for bubble formation
        \item Tighten connections as needed
    \end{itemize}
    \item \textbf{Inspect autosampler:}
    \begin{itemize}
        \item Clean carousel and boat holders
        \item Check positioning mechanism alignment
        \item Lubricate moving parts per manufacturer specifications
        \item Test carousel rotation (manual and automatic)
    \end{itemize}
\end{enumerate}

\subsection{Monthly Maintenance}

\begin{enumerate}[noitemsep]
    \item \textbf{Replace oxygen scrubber cartridge:}
    \begin{itemize}
        \item Locate moisture/oxygen scrubber in gas line
        \item Record installation date of current cartridge
        \item Replace per manufacturer schedule (typically 30 days)
        \item Install new cartridge with proper orientation
        \item Purge gas lines after replacement
    \end{itemize}
    \item \textbf{Inspect all O-rings and seals:}
    \begin{itemize}
        \item Check furnace door seals
        \item Inspect gas line O-rings
        \item Look for cracks, deformation, or hardening
        \item Replace worn seals
        \item Apply vacuum grease to new O-rings
    \end{itemize}
    \item \textbf{Clean or replace furnace tubes:}
    \begin{itemize}
        \item Remove tubes and inspect for damage
        \item Heavy carbon buildup: replace tube
        \item Light buildup: clean with ceramic brush
        \item Soak stubborn deposits in dilute nitric acid (10\%)
        \item Rinse thoroughly with deionized water
        \item Dry at 105°C before reinstallation
    \end{itemize}
    \item \textbf{Perform full system calibration:}
    \begin{itemize}
        \item Prepare fresh calibration standards
        \item Run complete calibration sequence (6 levels)
        \item Verify system performance meets specifications
        \item Document calibration in maintenance log
    \end{itemize}
    \item \textbf{Analyze certified reference material:}
    \begin{itemize}
        \item Use NIST or equivalent CRM
        \item Analyze in quintuplicate
        \item Calculate recovery and precision
        \item Must be within certified range
        \item Document in QC control charts
    \end{itemize}
    \item \textbf{Update maintenance log:}
    \begin{itemize}
        \item Record all maintenance performed
        \item Note any parts replaced (with part numbers)
        \item Document any issues observed
        \item Sign and date maintenance record
    \end{itemize}
\end{enumerate}

\subsection{Common Issues and Solutions}

\begin{longtable}{p{3.5cm}p{4cm}p{6.5cm}}
\toprule
\textbf{Problem} & \textbf{Possible Cause} & \textbf{Solution} \\
\midrule
\endfirsthead
\multicolumn{3}{c}%
{\tablename\ \thetable\ -- \textit{Continued from previous page}} \\
\toprule
\textbf{Problem} & \textbf{Possible Cause} & \textbf{Solution} \\
\midrule
\endhead
\midrule
\multicolumn{3}{r}{\textit{Continued on next page}} \\
\endfoot
\bottomrule
\endlastfoot
High blank values ($>$ 0.1 mg C) & Contaminated boats or combustion tube & Pre-combust boats at 900°C for 15 min; clean combustion tube; replace quartz wool \\
Poor peak shape (tailing) & Incomplete combustion or contamination & Increase combustion time to 5 min; raise temperature to 920°C; clean combustion tube \\
Low TIC recovery & Insufficient acid or incomplete reaction & Increase acid volume to 200--250 $\mu$L; extend heating time to 10 min; use higher acid concentration (30\%) \\
Vigorous reaction with high carbonate soils & Rapid CO$_2$ evolution & Reduce sample size to 5--10 mg; add acid in small increments (50 $\mu$L); wait between additions \\
Baseline drift ($>$ 5 ppm/hour) & Detector contamination or temperature instability & Clean NDIR detector cell; verify furnace temperature stability; check gas flow rates \\
Inconsistent replicates (RSD $>$ 15\%) & Sample heterogeneity or weighing errors & Improve grinding (finer particle size); increase mixing time; check balance calibration \\
Negative TOC values & TIC $>$ TC measurement error & Reanalyze both TC and TIC; verify calibrations; check sample weights; ensure no sample loss \\
Peak not detected & Sample too small or detector sensitivity low & Increase sample size; check detector settings; verify calibration range; clean detector \\
Carryover between samples & Incomplete purging or residue in tube & Extend purge time to 60 sec; increase furnace temperature; clean combustion tube more frequently \\
Software communication error & Cable disconnection or driver issue & Check USB/RS-232 connection; restart software; reboot computer; reinstall drivers if necessary \\
Furnace not heating & Heating element failure or power issue & Check power supply; verify temperature controller settings; test heating element continuity; replace if faulty \\
Acid not dispensing & Syringe blockage or empty reservoir & Refill acid reservoir; clean syringe needle; prime system; check for air bubbles in line \\
Autosampler malfunction & Mechanical jam or position error & Manually rotate carousel to identify obstruction; check boat placement; clean positioning sensors; recalibrate positions \\
Oxygen pressure drop & Leak or cylinder empty & Check cylinder pressure; perform leak test on all fittings; tighten connections; replace cylinder if low \\
Moisture in gas lines & Saturated desiccant or condensation & Replace desiccant in moisture trap; check for water in gas lines; install heated gas lines if necessary \\
\end{longtable}

\subsection{When to Call for Service}

Contact Shimadzu service technician if:

\begin{itemize}[noitemsep]
    \item Furnace temperatures cannot be maintained or controlled
    \item NDIR detector shows erratic or no response after cleaning
    \item Persistent baseline drift ($>$ 10 ppm/hour) despite troubleshooting
    \item Calibration cannot achieve $r^2 > 0.99$ with known good standards
    \item Autosampler positioning repeatedly fails
    \item Gas leaks cannot be eliminated by tightening fittings
    \item Software errors persist after reinstallation
    \item Any unusual sounds, smells, or visible damage to instrument
    \item Annual preventive maintenance is due
\end{itemize}

\section{Waste Disposal}

\subsection{Chemical Waste}

\begin{itemize}[noitemsep]
    \item \textbf{Phosphoric acid waste:}
    \begin{itemize}
        \item Collect in designated acid waste container
        \item Neutralize with solid sodium bicarbonate (NaHCO$_3$) before disposal
        \item Add bicarbonate slowly until pH 6--8 (test with pH paper)
        \item Dispose via institutional chemical waste program
        \item Label container: "Neutralized Phosphoric Acid Waste"
    \end{itemize}
    \item \textbf{Standard solutions:}
    \begin{itemize}
        \item Glucose/sucrose solutions: dilute and dispose to sanitary sewer (if permitted)
        \item Calcium carbonate: can be neutralized and disposed as solid waste
        \item Store expired standards separately, labeled with disposal date
    \end{itemize}
    \item \textbf{Cleaning solvents:}
    \begin{itemize}
        \item Methanol, acetone: collect in organic waste container
        \item Dispose per institutional hazardous waste procedures
    \end{itemize}
\end{itemize}

\subsection{Solid Waste}

\begin{itemize}[noitemsep]
    \item \textbf{Ceramic boats:}
    \begin{itemize}
        \item Reuse after pre-combustion cleaning (typical life: 50--100 uses)
        \item Discard damaged boats as non-hazardous solid waste
        \item If contaminated with heavy metals: dispose as hazardous waste
    \end{itemize}
    \item \textbf{Analyzed soil samples:}
    \begin{itemize}
        \item Non-contaminated soils: dispose as non-hazardous solid waste
        \item Known contaminated soils: return to project manager or dispose per waste characterization
    \end{itemize}
    \item \textbf{Quartz wool and furnace packing:}
    \begin{itemize}
        \item Non-hazardous solid waste
        \item Discard when discolored or compressed
    \end{itemize}
    \item \textbf{Desiccant:}
    \begin{itemize}
        \item Indicating silica gel: regenerate by heating at 200°C for 2 hours, or discard as non-hazardous waste
        \item Drierite: dispose as non-hazardous solid waste when saturated
    \end{itemize}
\end{itemize}

\subsection{Gas Cylinders}

\begin{itemize}[noitemsep]
    \item Return empty oxygen cylinders to supplier
    \item Do not dispose in regular waste
    \item Close cylinder valve before returning
    \item Mark cylinder as "EMPTY" with tag
\end{itemize}

\section{References}

\begin{enumerate}
    \item Shimadzu SSM-5000A Total Organic Carbon Analyzer User Manual, Rev. 3.0, Shimadzu Corporation, 2023
    \item ASTM D4129-13: Standard Test Method for Total and Organic Carbon in Water by High Temperature Oxidation and by Coulometric Detection
    \item ISO 10694:1995: Soil Quality -- Determination of Organic and Total Carbon after Dry Combustion (Elementary Analysis)
    \item Nelson, D.W. and Sommers, L.E. (1996) Total Carbon, Organic Carbon, and Organic Matter. In: Methods of Soil Analysis Part 3 -- Chemical Methods, SSSA Book Series 5, Soil Science Society of America, Madison, WI, pp. 961--1010
    \item USEPA Method 9060A: Total Organic Carbon, SW-846 Test Methods for Evaluating Solid Waste
    \item Loeppert, R.H. and Suarez, D.L. (1996) Carbonate and Gypsum. In: Methods of Soil Analysis Part 3 -- Chemical Methods, SSSA Book Series 5, pp. 437--474
\end{enumerate}

\section{Revision History}

\begin{table}[h]
\centering
\begin{tabular}{cccc}
\toprule
\textbf{Version} & \textbf{Date} & \textbf{Changes} & \textbf{Author} \\
\midrule
1.0 & October 2025 & Initial release with detailed startup, gas system, standard curve, and sample table procedures & \\
\bottomrule
\end{tabular}
\end{table}

\section{Approval}

\begin{table}[h]
\centering
\begin{tabular}{cccc}
\toprule
\textbf{Role} & \textbf{Name} & \textbf{Signature} & \textbf{Date} \\
\midrule
Author & & & \\[0.5cm]
Technical Reviewer & & & \\[0.5cm]
Quality Assurance Manager & & & \\[0.5cm]
Laboratory Director & & & \\[0.5cm]
\bottomrule
\end{tabular}
\end{table}

\end{document}
